% reportmd が生成した LaTeX のプリアンブル（設定ファイルの内容を反映）
\usepackage[b5paper,top=25mm,bottom=25mm,left=3cm,right=25mm]{geometry}
\usepackage{xeCJK}
% 本文の太字（**…**）。太字の書体がない本文フォント（IPAex明朝など）でも太く見えるよう、疑似太字を使う
\setmainfont{IPAexMincho}[AutoFakeBold=2.5]
\setCJKmainfont{IPAexMincho}[AutoFakeBold=2.5]
\setmonofont{Noto Sans Mono CJK JP}
\setCJKmonofont{Noto Sans Mono CJK JP}
\usepackage{indentfirst}
\newlength{\reportmdchar}
\setlength{\reportmdchar}{11pt}
\setlength{\parindent}{0pt}
\linespread{1.25}
\makeatletter
\renewcommand\normalsize{\@setfontsize\normalsize{11pt}{13.2pt}}
\makeatother
\pagestyle{empty}
% --- 見出し（レベル1〜3：section, subsection, subsubsection） ---
\newfontfamily\reportmdheadlatin{IPAexGothic}
\newCJKfontfamily\reportmdheadcjk{IPAexGothic}
\makeatletter
\newcommand{\reportmdheading}[2]{\reportmdheadlatin\reportmdheadcjk\fontsize{#1}{#2}\selectfont}
\renewcommand{\@seccntformat}[1]{\csname the#1\endcsname\hspace{1em}}
\renewcommand{\section}{\@startsection{section}{1}{0pt}{16.5pt}{8.25pt}{\reportmdheading{14pt}{16.8pt}}}
\renewcommand{\subsection}{\@startsection{subsection}{2}{5.5pt}{16.5pt}{8.25pt}{\reportmdheading{12pt}{14.4pt}}}
\renewcommand{\subsubsection}{\@startsection{subsubsection}{3}{11pt}{16.5pt}{8.25pt}{\reportmdheading{10.5pt}{12.6pt}}}
\makeatother
% --- 表（filter.lua が tabular に組み、ここで定義した命令で書式を決める） ---
\usepackage{array}
\usepackage{calc}
\usepackage{tabularx}
\newenvironment{reportmdtablebox}{\begin{flushleft}}{\end{flushleft}}
\newcommand{\reportmdtablestyle}{\linespread{1}\fontsize{9pt}{10.8pt}\selectfont\setlength{\tabcolsep}{6pt}\renewcommand{\arraystretch}{1.3}}
% 見出し行の太字。太字の書体がない本文フォント（IPAex明朝など）でも太く見えるよう、疑似太字を使う
\newfontfamily\reportmdtableheadlatin{IPAexMincho}[AutoFakeBold=2.5]
\newCJKfontfamily\reportmdtableheadcjk{IPAexMincho}[AutoFakeBold=2.5]
\newcommand{\reportmdtablehead}[1]{{\reportmdtableheadlatin\reportmdtableheadcjk\bfseries #1}}
% 横の罫線（使わない線は、空の命令にする）
\newcommand{\reportmdruletop}{\noalign{\hrule height 0.8pt}}
\newcommand{\reportmdruleheader}{\noalign{\hrule height 1.2pt}}
\newcommand{\reportmdrulebottom}{\noalign{\hrule height 0.8pt}}
\newcommand{\reportmdrulerows}{}
% 縦の罫線（列の間 vrule と、表の左右の端 vside）
\newcommand{\reportmdvrule}{}
\newcommand{\reportmdvside}{}
